\rtitle{Evolution of Traffic Flow Visualization in Civil Engineering Education}
\rsec{Introduction}
The visualization of traffic flow theory and highway capacity analysis has undergone a remarkable transformation in civil engineering education. From hand-drawn \clue{time-space diagrams} and \clue{manual traffic count sheets} to sophisticated \clue{microsimulation animations} and \clue{real-time data dashboards}, the tools available to students and practitioners have evolved dramatically. This report examines how \clue{textbooks, laboratory exercises, and student projects} have represented congestion patterns, level of service concepts, and intersection operations through different visual media, tracing the progression from classical analytical methods to modern computational approaches.

\rsec{Classical Foundations: Time-Space Diagrams and Fundamental Relationships}
\rnote{Every concept the request names (time-space diagrams, count sheets, microsimulation, dashboards, LOS) was searched in the first turn; the time-space figure used here was inspected in turn 3, right after its Wikipedia page was opened.}
The foundation of traffic flow visualization rests on time-space diagrams, which plot vehicle position against time to reveal movement patterns and congestion propagation. These diagrams remain essential teaching tools because they make abstract concepts tangible.

\rtrace{turn 2 \texttt{open(url)} \(\rightarrow\) turn 3 \texttt{open(id=page\_img\_32)}}{Source page: \emph{Traffic flow - Wikipedia}; caption on the page: ``Figure 1. Time Space diagram''\par\textbf{Why it was used.} The figure marks exactly the quantities the paragraph explains: the gaps between trajectories (\(h_3\), \(s_3\)) and the slope \(v\) of one trajectory.}
\rimagebox{figures/showcase/traffic/detail/page_img_32__Time_Space_Diag_Figure_1__e01e89d919.jpg}{Time-space diagram showing vehicle trajectories with labeled spacing and headway measurements}{\rbox{0.4}{0.43}{0.75}{0.53}{}\rbox{0.35}{0.27}{0.43}{0.5}{}\rbox{0.72}{0.71}{0.8}{0.82}{}}
The time-space diagram above illustrates how individual vehicle trajectories can be analyzed to extract key parameters. The vertical spacing between trajectories at a given time represents the time headway, while horizontal spacing at a given position indicates the space headway. The slope of each trajectory represents vehicle speed, and the curvature reveals acceleration and deceleration patterns \rcite{1}. This visualization technique allows students to understand shock wave propagation, queue formation, and the fundamental relationship between flow, density, and speed.

The \clue{fundamental diagram of traffic flow} provides the mathematical backbone for these visualizations. The relationship between traffic flow (Q), density (D), and velocity (V) is expressed as \fact{Q = D\(\cdot\)V}, with the diagram showing distinct operating regimes: free flow, bound flow, and congestion \rcite{9}.

\rtrace{turn 11 \texttt{open(url)} \(\rightarrow\) turn 12 \texttt{open(id=page\_img\_183)}}{Source page: \emph{Fundamental diagram of traffic flow - Wikipedia}; caption on the page: ``Fundamental diagram of traffic flow''\par\textbf{Why it was used.} Chosen for the boxed identity at the top right, which the text quotes as Q = D\(\cdot\)V, and for the three labelled regimes in the top panel.}
\rimagebox{figures/showcase/traffic/detail/page_img_183__500px-Fundamental_Diagram__447a6ca776.jpg}{Fundamental diagram of traffic flow showing free flow, bound flow, and congestion regimes with velocity-density relationships}{\rbox{0.71}{0.905}{0.92}{0.966}{}\rbox{0.12}{0.7}{0.65}{0.865}{regimes}}
The fundamental diagram reveals critical insights for education. The free flow regime shows vehicles traveling at or near free-flow speed with low density. As density increases, vehicles enter the bound flow regime where speeds begin to decrease. At maximum density (D\_max), congestion occurs and flow drops to zero. The triangular flow-density curve, widely taught in academic settings, consists of a free-flow branch and a congested branch with negative slope, with the apex representing roadway capacity \rcite{9}.

\rtrace{turn 2 \texttt{open(url)} \(\rightarrow\) turn 3 \texttt{open(id=page\_img\_34)}}{Source page: \emph{Traffic flow - Wikipedia}; caption on the page: ``Figure 3. Flow Density relationship''\par\textbf{Why it was used.} Supplies the four quantities the paragraph defines: critical density, jam density, wave speed and free-flow speed.}
\rimagebox{figures/showcase/traffic/detail/page_img_34__960px-Flow_Density_Relationship__57ffd32caf.jpg}{Flow-density relationship showing critical density, jam density, free flow speed, and wave speed}{\rbox{0.14}{0.18}{0.44}{0.24}{}\rbox{0.68}{0.18}{0.93}{0.24}{}\rbox{0.42}{0.49}{0.68}{0.565}{}\rbox{0.11}{0.41}{0.43}{0.48}{}}
This flow-density relationship diagram demonstrates how wave speed (\(\omega\)) relates to the slope of the congested branch, while free flow speed (V\_f) relates to the slope of the free-flow branch. Students learn that the critical density (k\_c) represents the density at maximum flow, and jam density (k\_j) represents the density when vehicles are bumper-to-bumper \rcite{1}.

\rsec{Level of Service and Capacity Analysis Visualization}
\rnote{LOS was the hardest concept to source: it appears in search queries in turns 1, 4, 6, 9 and 10. The Wikipedia LOS page that finally anchors the A--F scale was opened in turn 7.}
The Highway Capacity Manual (HCM) introduced \clue{Level of Service (LOS)} as a qualitative measure to relate traffic quality, categorizing conditions \fact{from A (free flow) to F (breakdown flow)} \rcite{4}. Visual representations of LOS have evolved from simple letter grades to sophisticated multi-dimensional analyses.

\rtrace{turn 2 \texttt{open(url)} \(\rightarrow\) turn 3 \texttt{open(id=page\_img\_35)}}{Source page: \emph{Traffic flow - Wikipedia}; caption on the page: ``Figure 4. Relationship between flow ( q ), density ( k ), and speed ( v )''}
\rimagebox{figures/showcase/traffic/detail/page_img_35__960px-Relationship_between_q_k_and_v__a0f722dca1.jpg}{Relationship between flow (q), density (k), and speed (v) showing interconnected traffic stream variables}{}
This interconnected diagram shows how flow, density, and speed relate across different traffic states. The triangular relationship between these variables forms the basis for capacity analysis, with the peak representing maximum throughput \rcite{1}. Students use these relationships to understand why operating near capacity is risky and how small increases in density can lead to disproportionate decreases in speed and flow.

\rtrace{turn 2 \texttt{open(url)} \(\rightarrow\) turn 3 \texttt{open(id=page\_img\_45)}}{Source page: \emph{Traffic flow - Wikipedia}; caption on the page: ``Speed-Flow Diagram for Typical Roadway''}
\rimagebox{figures/showcase/traffic/detail/page_img_45__Speed-Flow_Diagram_for_Typical_Roadway__ae84e8eeaf.jpg}{Speed-flow diagram for typical roadway showing the characteristic curve shape}{}
The speed-flow diagram illustrates the characteristic relationship where speed decreases as flow increases, with the curve peaking at capacity. The stable regime (left of peak) allows drivers to adjust speeds freely, while the unstable regime (right of peak) leads to breakdown and stop-and-go conditions. This visualization helps students understand the transition from free flow to congested conditions \rcite{1}.

\rtrace{turn 2 \texttt{open(url)} \(\rightarrow\) turn 3 \texttt{open(id=page\_img\_46)}}{Source page: \emph{Traffic flow - Wikipedia}; caption on the page: ``Figure 18''}
\rimagebox{figures/showcase/traffic/detail/page_img_46__960px-Wikipedia_TrafficBottlenecks_regular.svg__828e2ce138.jpg}{Figure showing bottleneck effects on traffic flow with shock wave propagation}{}
This figure demonstrates how bottlenecks create shock waves that propagate upstream, reducing effective capacity. The Federal Highway Authority \fact{attributes 40\% of congestion to bottlenecks}, making this visualization essential for understanding real-world traffic dynamics \rcite{1}.

\rsec{Microsimulation and Dynamic Visualization}
The advent of powerful computing has enabled microsimulation, where individual vehicle behaviors are modeled and animated. This represents a paradigm shift from static diagrams to dynamic, interactive visualizations.

\rtrace{turn 2 \texttt{open(url)} \(\rightarrow\) turn 3 \texttt{open(id=page\_img\_56)}}{Source page: \emph{Traffic Simulation Lab - Transportation Engineering and Science Program}; caption on the page: ``Traffic Simulation Lab''\par\textbf{Why it was used.} The lab's own photo, used to ground the claim about dedicated simulation laboratories.}
\rimagebox{figures/showcase/traffic/detail/page_img_56__traffic-simulation-lab-1-768x392__aab037cf4c.jpg}{Traffic Simulation Lab with students working at computers with large display screen}{}
University transportation engineering programs now maintain dedicated \clue{simulation laboratories}. The University of Tennessee's Traffic Simulation Lab, for example, houses high-performance computers with software including PTV VISSIM, Aimsun, and Synchro, \fact{supporting up to 50 students simultaneously} \rcite{2}. These facilities enable students to experience variable transportation simulation software and tools in a collaborative environment.

\rtrace{turn 2 \texttt{open(url)} \(\rightarrow\) turn 3 \texttt{open(id=page\_img\_57)}}{Source page: \emph{Traffic Simulation Lab - Transportation Engineering and Science Program}; caption on the page: ``Students Working in the Traffic Simulation Lab''}
\rimagebox{figures/showcase/traffic/detail/page_img_57__traffic-sumulation-lab-gallery-1-300x224__7a846fc196.jpg}{Students working in the Traffic Simulation Lab with computers and large projection screen}{}
Microsimulation software like PTV VISSIM enables traffic planners and engineers to ``plan, test and justify mobility measures before they are implemented'' through ``realistic microscopic and multimodal simulation, powerful scenario evaluation, and convincing visual outputs'' \rcite{7}. The software digitally reproduces traffic patterns of all road users on a microscopic scale, using science-based motion and interaction models for pedestrians, vehicles, micromobility, and automated vehicles.

\rtrace{turn 2 \texttt{open(url)} \(\rightarrow\) turn 3 \texttt{open(id=page\_img\_58)}}{Source page: \emph{Traffic Simulation Lab - Transportation Engineering and Science Program}; caption on the page: ``Students Working in the Traffic Simulation Lab''}
\rimagebox{figures/showcase/traffic/detail/page_img_58__traffic-sumulation-lab-gallery-2-300x224__aab1d8ff85.jpg}{Students collaborating in the Traffic Simulation Lab with multiple workstations}{}
\clue{Student projects} in these labs typically involve building network models, calibrating parameters against field data, and running scenarios to evaluate design alternatives. The visual output---animated vehicle movements, queue formations, and signal operations---provides immediate feedback that static diagrams cannot match.

\rsec{Real-Time Data Dashboards and Modern Visualization}
\rnote{The three dashboard figures come from one vendor page opened in turn 7 and were inspected together in turn 8; the report uses them to illustrate KPIs, not as evidence of classroom use.}
Contemporary traffic engineering increasingly relies on \clue{real-time data dashboards that aggregate} information from multiple sources into intuitive visual interfaces. These tools have become essential for both professional practice and educational applications.

\rtrace{turn 7 \texttt{open(url)} \(\rightarrow\) turn 8 \texttt{open(id=page\_img\_134)}}{Source page: \emph{Traffic Management Dashboards: Improving Flow, Safety, and Urban Mobility}; caption on the page: ``traffic management dashboard example''}
\rimagebox{figures/showcase/traffic/detail/page_img_134__traffic-management-dashboard__ab6e799d75.jpg}{Traffic management dashboard showing total incidents, congestion levels, average speed, and vehicles per hour with trend lines}{}
Modern dashboards present key performance indicators including total incidents, average congestion level, average speed, and vehicles per hour, with year-over-year comparisons and trend lines \rcite{6}. The dashboard format allows operators to monitor conditions, respond to incidents, and optimize traffic flow across entire regions.

\rtrace{turn 7 \texttt{open(url)} \(\rightarrow\) turn 8 \texttt{open(id=page\_img\_135)}}{Source page: \emph{Traffic Management Dashboards: Improving Flow, Safety, and Urban Mobility}; caption on the page: ``traffic congestion dashboard''}
\rimagebox{figures/showcase/traffic/detail/page_img_135__traffic-congestion-dashboard__999331f749.jpg}{Traffic congestion dashboard with detailed metrics and visual indicators}{}
These dashboards integrate data from traffic sensors, cameras, GPS devices, weather systems, and public transit feeds, transforming raw data into intuitive charts, maps, and performance indicators \rcite{6}. Common KPIs include average vehicle speed, traffic volume, travel time index, and incident response time.

\rtrace{turn 7 \texttt{open(url)} \(\rightarrow\) turn 8 \texttt{open(id=page\_img\_136)}}{Source page: \emph{Traffic Management Dashboards: Improving Flow, Safety, and Urban Mobility}; caption on the page: ``traffic incident dashboard''\par\textbf{Why it was used.} Its incident list and sensor panel are what the text calls a \emph{single pane of glass}.}
\rimagebox{figures/showcase/traffic/detail/page_img_136__traffic-incident-dashboard__4c2bed7726.jpg}{Traffic incident dashboard showing incident summary, signal performance, and roadway throughput}{\rbox{0.025}{0.05}{0.465}{0.49}{incidents}\rbox{0.475}{0.05}{0.98}{0.49}{sensors}}
The incident dashboard format demonstrates how multiple data streams are consolidated into a single pane of glass, enabling traffic control centers to detect congestion hotspots, identify accidents quickly, coordinate emergency response, and adjust traffic signals dynamically \rcite{6}.

\rsec{Space-Time Diagrams in Modern Context}
\rnote{Turn 10 added the query \emph{traffic flow theory fundamental diagram speed flow density textbook}; the network-link diagram and the MFD below were inspected in turn 12 from the page opened in turn 11.}
\clue{Space-time diagrams} continue to serve as fundamental teaching tools, now often generated computationally from simulation outputs or field data.

\rtrace{turn 11 \texttt{open(url)} \(\rightarrow\) turn 12 \texttt{open(id=page\_img\_186)}}{Source page: \emph{Fundamental diagram of traffic flow - Wikipedia}; caption on the page: ``Figure 2: Space-time diagram for the ith link in a traffic flow network''\par\textbf{Why it was used.} Its rectangles A and B are the two regions the paragraph refers to.}
\rimagebox{figures/showcase/traffic/detail/page_img_186__500px-Time_Space_Diag_Figure_2__e263c2d408.jpg}{Space-time diagram for the ith link in a traffic flow network showing vehicle trajectories}{\rbox{0.385}{0.26}{0.465}{0.84}{}\rbox{0.18}{0.49}{0.84}{0.595}{}}
This space-time diagram for a network link shows how vehicle trajectories can be analyzed to understand queue formation and discharge. The rectangular regions A and B represent different traffic states, with the slope of trajectories indicating speed and the vertical spacing indicating headway \rcite{9}.

\rtrace{turn 11 \texttt{open(url)} \(\rightarrow\) turn 12 \texttt{open(id=page\_img\_185)}}{Source page: \emph{Fundamental diagram of traffic flow - Wikipedia}; caption on the page: ``Figure 1: Sample traffic flow macroscopic fundamental diagram''\par\textbf{Why it was used.} The labelled \emph{sweet spot} is the operating point the paragraph names.}
\rimagebox{figures/showcase/traffic/detail/page_img_185__500px-Sample_Traffic_Flow_Macroscopic_Fundamental_Diagram__b09c400342.jpg}{Sample traffic flow macroscopic fundamental diagram showing aggregate network behavior}{\rbox{0.29}{0.12}{0.45}{0.19}{}}
Modern \clue{macroscopic fundamental diagrams (MFDs)} extend the fundamental diagram concept to network-wide analysis, showing the relationship between average network flow and average network density. The ``sweet spot'' region represents optimal network operation, while the right side indicates congestion propagation \rcite{9}.

\rsec{Educational Integration and Pedagogical Evolution}
The evolution of visualization tools has fundamentally changed how traffic flow concepts are taught. Early curricula relied heavily on hand calculations and static diagrams. Modern programs integrate multiple visualization approaches:

\begin{renum}
\item \textbf{Classical analytical methods}: Time-space diagrams and fundamental diagrams teach the underlying physics and mathematics
\item \textbf{Microsimulation}: Dynamic animations provide intuitive understanding of complex interactions
\item \textbf{Real-time dashboards}: Connect classroom learning to operational practice
\item \textbf{Interactive web-based tools}: Enable remote learning and collaborative analysis
\end{renum}
The Highway Capacity Manual itself has evolved to incorporate these visualization approaches, with the 2010 edition incorporating multimodal analysis tools and \fact{the 2022 7th Edition introducing major methodological changes} \rcite{3}. Professional development courses now teach students to use \clue{Highway Capacity Software (HCS)} to perform HCM-based capacity analyses, with visual outputs supporting interpretation \rcite{3}.

\rsec{Conclusion}
The visualization of traffic flow theory has evolved from simple hand-drawn diagrams to sophisticated multi-modal visual environments. Time-space diagrams and fundamental diagrams remain essential for teaching core concepts, while microsimulation animations provide dynamic, intuitive understanding of complex interactions. Real-time dashboards bridge the gap between academic learning and professional practice. Together, these visualization tools enable civil engineering students to develop both theoretical understanding and practical skills for analyzing and managing traffic systems. The progression from static to dynamic, from analytical to computational, and from isolated to integrated visualization reflects broader trends in engineering education toward experiential, technology-enhanced learning.

\rsec{References}
\rcite{1} Traffic flow - Wikipedia\rurl{https://en.wikipedia.org/wiki/Traffic_flow}

\rcite{2} Traffic Simulation Lab - Transportation Engineering and Science Program\rurl{https://tesp.utk.edu/facilities/traffic-simulation-lab/}

\rcite{3} Highway Capacity Analysis - McTrans Center\rurl{https://mctrans.ce.ufl.edu/training/hca/}

\rcite{4} Level of service (transportation) - Wikipedia\rurl{https://en.wikipedia.org/wiki/Level_of_service_(transportation)}

\rcite{5} Read ``Traffic Capacity Level of Service: Adaptations and Usage'' at NAP.edu\rurl{https://www.nationalacademies.org/read/29143/chapter/4}

\rcite{6} Traffic Management Dashboards: Improving Flow, Safety, and Urban Mobility\rurl{https://www.inetsoft.com/info/traffic-management-dashboard/}

\rcite{7} PTV Vissim\rurl{https://www.ptvgroup.com/en-us/products/ptv-vissim}

\rcite{8} Traffic Engineering - ITE-A Community of Transportation Professionals\rurl{https://www.ite.org/technical-resources/topics/traffic-engineering/}

\rcite{9} Fundamental diagram of traffic flow - Wikipedia\rurl{https://en.wikipedia.org/wiki/Fundamental_diagram_of_traffic_flow}

\rcite{10} 5.2: Traffic Flow - Engineering LibreTexts\rurl{https://eng.libretexts.org/Bookshelves/Civil_Engineering/Fundamentals_of_Transportation/05\%3A_Traffic/5.02\%3A_Traffic_Flow}

\rgate{\item Submission 1 (turn 15): {\color{rframe}\textbf{WARN}} --- 8,586 chars, 11 images, 10 sources; \texttt{unexamined\_image}\,\(\times\)2, \texttt{nonlocal\_or\_missing\_destination}\,\(\times\)2
\item Submission 2 (turn 17): {\color{tok}\textbf{OK}} --- 8,861 chars, 14 images, 10 sources}
